\chapter*{Phụ Lục III. Bốn Tuần Tập Sống Chắt Chiu}
\addcontentsline{toc}{chapter}{Phụ Lục III. Bốn Tuần Tập Sống Chắt Chiu}
\markboth{Bốn Tuần Tập Sống Chắt Chiu}{Bốn Tuần Tập Sống Chắt Chiu}

\section*{Tuần 1. Nhìn Thẳng Vào Dòng Tiền}
\begin{truyen}{Bảy Ngày Không Trốn Tránh}{Seven Days of Honest Observation}
Trong bảy ngày đầu, hãy ghi toàn bộ chi tiêu, kể cả những khoản rất nhỏ. Không cần đẹp, chỉ cần thật. Khi chưa nhìn ra dòng tiền, chưa thể sửa dòng tiền.\
\textit{(During the first seven days, record all spending, including the smallest expenses. It need not be beautiful, only honest. Until you see the flow of money, you cannot correct it.)}

Mỗi tối, khoanh tròn ba khoản tiêu mà nếu bình tĩnh hơn, bạn đã có thể tránh. Chưa cần sửa hết. Chỉ cần nhìn ra.\
\textit{(Each evening, circle three expenses that could have been avoided with more calmness. There is no need to fix everything yet. First, simply see it.)}
\end{truyen}

\section*{Tuần 2. Chia Tiền Thành Bốn Ngăn}
\begin{truyen}{Đừng Để Tất Cả Tiền Nằm Chung Một Cảm Giác}{Do Not Leave All Money in One Emotional Pile}
Tuần thứ hai, hãy chia thu nhập thành bốn phần: sinh hoạt thiết yếu, học hành và gia đình, quỹ dự phòng, khoản để dành lâu hơn. Dù ít cũng phải chia.\
\textit{(In the second week, divide income into four parts: essential living, family and education, emergency fund, and longer-term saving. Even if the amounts are small, divide them.)}

Việc chia này không phải để làm màu. Nó để tiền có ranh giới, còn lòng mình bớt tùy tiện.\
\textit{(This division is not for appearance. It gives money boundaries and makes the heart less careless.)}
\end{truyen}

\section*{Tuần 3. Cắt Một Thói Quen Lãng Phí}
\begin{truyen}{Không Cần Sửa Tất Cả Cùng Lúc}{You Do Not Need to Fix Everything at Once}
Chọn một thói quen đang rò tiền nhiều nhất: cà phê ngoài, mua đồ giảm giá vô thức, gọi món vì lười nấu, hay tiêu vì buồn. Chỉ chọn một.\
\textit{(Choose the single habit that leaks the most money: coffee outside, mindless sale purchases, ordering food out of laziness, or spending because of sadness. Choose only one.)}

Trong bảy ngày, tập thay thế nó bằng một cách khác nhẹ hơn. Khi một lỗ rò được chặn lại, lòng tin vào chính mình cũng mạnh lên.\
\textit{(For seven days, practice replacing it with a gentler alternative. When one leak is stopped, trust in oneself also grows stronger.)}
\end{truyen}

\section*{Tuần 4. Dạy Lại Cả Nhà Một Nếp Mới}
\begin{truyen}{Tiết Kiệm Không Phải Việc Của Một Người}{Saving Is Not One Person’s Burden Alone}
Tuần cuối, hãy nói chuyện thẳng thắn với gia đình về bốn điều: điều gì đang làm nhà mình tốn tiền, điều gì thật sự quan trọng, quỹ nào cần giữ, và con cái nên học điều gì từ cha mẹ.\
\textit{(In the final week, speak honestly with the family about four things: what is costing the household money, what truly matters, which fund must be protected, and what children should learn from the parents.)}

Không cần nói nặng nề. Chỉ cần rõ ràng. Một mái nhà bền không thể chỉ dựng bằng sự âm thầm chịu đựng của một người.\
\textit{(There is no need to speak harshly. Clarity is enough. A durable home cannot be built solely on the silent endurance of one person.)}
\end{truyen}

\section*{Mười Câu Tự Hỏi Mỗi Cuối Tuần}
\begin{truyen}{Những Câu Hỏi Để Không Quay Lại Lối Cũ}{Questions That Prevent a Return to Old Habits}
1. Tuần này tôi đã tiêu vì điều gì là chính: nhu cầu hay cảm xúc?\
\textit{(This week, did I spend mainly for needs or emotions?)}

2. Tôi còn nhớ đồng tiền mình kiếm ra đã vất vả thế nào không?\
\textit{(Do I still remember how hard-earned my money is?)}

3. Khoản nào là lỗ rò lớn nhất của tuần này?\
\textit{(What was the biggest leak this week?)}

4. Tôi có để dành được dù chỉ một ít không?\
\textit{(Did I manage to save even a small amount?)}

5. Tôi có lần nào tiêu vì sĩ diện không?\
\textit{(Did I spend out of pride at any point?)}

6. Tôi có dạy con được điều gì bằng chính hành động của mình không?\
\textit{(Did I teach my children anything through my own actions?)}

7. Tôi có dùng nợ để che thói quen xấu không?\
\textit{(Did I use debt to cover a bad habit?)}

8. Tôi có biết ơn điều mình đang có không?\
\textit{(Was I grateful for what I already have?)}

9. Tôi đã làm gì để giữ ngày mai của gia đình?\
\textit{(What did I do to protect my family’s tomorrow?)}

10. Tuần tới tôi muốn sửa đúng một điều gì?\
\textit{(What one thing do I want to correct next week?)}
\end{truyen}

\begin{baihoc}
Thực hành đều quan trọng hơn hiểu nhiều. Chỉ cần mỗi tuần sửa được một lỗ rò, sau vài tháng đời sống đã khác.
\textit{(Consistent practice matters more than understanding many ideas. If one leak is repaired each week, life changes noticeably after a few months.)}
\end{baihoc}

\begin{ghinhoanh}
Practice beats intention.

Weekly honesty builds long-term stability.
\end{ghinhoanh}

\ngancach